\statement{A.82}{Proved}{The actual reflection augmentation module is projective at every prime and free at the dyadic prime}
Let \(E=F\rtimes C_2\), where \(F=\widehat F_2\), be
one of the actual pants extensions already considered: the pure
negative extension of C.28, or the regular extension of C.43.
Let \(c_1,\ldots,c_n\) represent its involution classes, put
\(C_i=\langle c_i\rangle\), and let \(\varepsilon\)
be its orientation character. Here \(n=3\) in the first case
and \(n=1\) in the second. For every prime p put
\[
 \Lambda_G=\mathbb Z_p[[G]],\qquad
 P_p=\bigoplus_{i=1}^n\mathbb Z_p[[E/C_i]],\qquad
 R_p=\ker(P_p\longrightarrow\mathbb Z_p),
 \tag{A.82.1}
\]
where the last map sums the augmentations. These are actual
compact modules over the full group E. For every finite discrete
E-module A, the restriction maps satisfy
\[
 \begin{aligned}
 H^1(E,A)&\longrightarrow\bigoplus_{i=1}^n H^1(C_i,A)
                    &&\text{is onto},\\
 H^2(E,A)&\longrightarrow\bigoplus_{i=1}^n H^2(C_i,A)
                    &&\text{is injective}.
 \end{aligned}
 \tag{A.82.2}
\]
Consequently \(R_p\) is finitely generated projective for
every p, and
\[
 R_p\oplus(R_p\otimes\varepsilon)\simeq\Lambda_E^{\,n+1},
 \qquad
 R_2\simeq\mathbb Z_2[[E]]^{(n+1)/2}.
 \tag{A.82.3}
\]
In particular the two dyadic ranks are two and one, respectively.
This does not produce a finite graph of groups or a generating
marked reflection system.

\begin{proof}
C.28 supplies the three classes in the pure negative case.
For the regular case we use the complete C.43 hypotheses:
\(F=\widehat F(x,y)\) is the full free subgroup of index
two, the odd involution interchanges the classes of x and y
with exponent \(-1\), and it preserves the class of
\(z=(xy)^{-1}\) with exponent \(-1\). C.43 proves
that there is exactly one F-conjugacy class of sections
and that its centralizer in F is trivial. Since F is
torsion-free, these are exactly the involution classes
in E, and in both cases \(C_E(c_i)=C_i\).
The index-two free subgroup and these self-centralizers
make E real projective, by the exact criterion in
\citet[Theorem~2.1]{A_PalQuickFormal2025}.

\emph{Prescribed involution classes are permitted.}\par
We use the allowed-class formulation of real projectivity in
\citet[Definition~1.1(4), Remark~1.2(b), and Proposition~2.3]{AHaran1990}.
We recall why the finite strengthening introduces no recognition
hypothesis. Given a finite epimorphism \(B\to Q\) and
a conjugacy-closed set \(\mathcal C\) of involutions in B,
form a finite-rank free profinite product D with one \(C_2\)
factor mapping to a representative of each class in
\(\mathcal C\), and enough free generators to map onto B.
The full free-product torsion theorem puts every involution
of D over \(\mathcal C\); see
\citet[Theorem~A']{A_HR1985}.

Let S be the preimage of \(B\setminus(\mathcal C\cup\{1\})\).
The compact set \(\{s^2:s\in S\}\) avoids one. Choose
an open normal U in \(\ker(D\to B)\) disjoint from that
set. Then \(D/U\to B\) lifts every element of
\(\mathcal C\) to an involution and has no involution
above a forbidden element of B. If the original finite
problem for E has nonidentity involution heads lying under
\(\mathcal C\), real projectivity solves the composite
problem through \(D/U\). Its projection to B sends all
involutions into \(\mathcal C\). This gives an ordinary
solution with the prescribed allowed classes; it does not
assert a proper solution.

\emph{Refine the finite action head before prescribing local
cocycles.}\par
Fix a finite E-module A and a finite action quotient
\(q_0:E\to Q_0\). There is a finer finite quotient
\(q:E\to Q\) in which all \(q(c_i)\) are nonidentity
and pairwise nonconjugate, and such that
\[
 \operatorname{im}\bigl(C_Q(q(c_i))\longrightarrow Q_0\bigr)
                         \subseteq q_0(C_i)
                 \quad(1\le i\le n).
 \tag{A.82.4}
\]
Nonidentity follows by retaining the orientation quotient.
Distinct compact conjugacy classes remain distinct in a
sufficiently fine finite quotient. For the centralizers,
as N ranges over open normal subgroups contained in
\(\ker q_0\), the closed sets
\[
                  X_{i,N}=\{g\in E:[g,c_i]\in N\}
\]
decrease to \(C_i\). If none had its \(q_0\)-image
inside \(q_0(C_i)\), compactness would give an element
of \(C_E(c_i)\) with a forbidden \(q_0\)-image.
This is impossible. A common refinement gives (A.82.4).
Only the projected centralizers are controlled here;
centralizers in arbitrary finite quotients need not be cyclic.

Choose arbitrary local classes represented by
\(a_i\in A\) with \((1+c_i)a_i=0\). In
\(B=A\rtimes Q\), each
\[
                            v_i=(a_i,q(c_i))
\]
is an involution. Use the union of their B-conjugacy classes
as the allowed set. The preceding finite problem has a
solution \(\psi:E\to B\), and we can write
\(\psi(g)=(d(g),q(g))\) for a continuous cocycle d.
The pairwise nonconjugate heads force \(\psi(c_i)\)
to be conjugate to \(v_i\), rather than another
\(v_j\). A conjugator \((a,k)\) has
\(k\in C_Q(q(c_i))\), and its formula is
\[
                         d(c_i)=ka_i+(1-c_i)a.
 \tag{A.82.5}
\]
By (A.82.4), k acts on A as one or as \(c_i\).
Either action leaves \([a_i]\in H^1(C_i,A)\)
unchanged. Thus d has all prescribed local classes,
proving the first assertion of (A.82.2). The centralizer
refinement is essential: allowed conjugacy classes alone
would not give the displayed local cohomology classes.

\emph{Degree two and compact projectivity.}\par
A continuous cocycle representing a class in \(H^2(E,A)\)
factors through a finite quotient through which the action
also factors. Refine that quotient to retain all involution
heads. If its restrictions to the \(C_i\) vanish, the
corresponding finite extension has involution lifts of all
these heads. A solution of the real embedding problem
splits its pullback to E. This proves the degree-two
injectivity in (A.82.2).

Continuous Ext and Shapiro applied to (A.82.1) give
\[
 \begin{gathered}
 H^1(E,A)\longrightarrow\bigoplus_iH^1(C_i,A)
       \longrightarrow\operatorname{Ext}^1_E(R_p,A)\\
       \longrightarrow H^2(E,A)
       \longrightarrow\bigoplus_iH^2(C_i,A).
 \end{gathered}
 \tag{A.82.6}
\]
For every finite p-primary A, the middle Ext group is
therefore zero. This is the compact projectivity criterion,
not merely a calculation with trivial coefficients. Indeed,
every map from \(R_p\) into a finite module quotient
lifts through a finite epimorphism by the vanishing of its
pullback extension class. The restriction computation below
makes \(R_p\) finitely generated. The finite lifting sets
are consequently finite, and compactness over the finite
quotients of any compact module epimorphism supplies a
continuous lift. Hence \(R_p\) is projective.

\emph{Split the induction counit, including at two.}\par
Each \(E/C_i\) restricts to the regular F-set, whence
\[
 \operatorname{Res}_F R_p
       \simeq\Lambda_F^{\,n-1}\oplus I_F
       \simeq\Lambda_F^{\,n+1}.
 \tag{A.82.7}
\]
The last equality uses the rank-two free augmentation ideal.
Induction and the tensor identity give the exact sequence
\[
 0\longrightarrow R_p\otimes\varepsilon
  \longrightarrow\operatorname{Ind}_F^E\operatorname{Res}_F R_p
  \longrightarrow R_p\longrightarrow0,
 \qquad
 \operatorname{Ind}_F^E\operatorname{Res}_F R_p
       \simeq\Lambda_E^{\,n+1}.
 \tag{A.82.8}
\]
The last map is the induction counit, and its kernel comes
from the augmentation ideal of \(\mathbb Z_p[E/F]\).
The sequence splits by the projectivity just proved.
At \(p=2\) this is not a splitting of
\(\mathbb Z_2[C_2]\) into sign eigenspaces.

Modulo two the twist disappears. Thus, with
\(\overline R=R_2/2R_2\),
\[
             \overline R\oplus\overline R
                     \simeq\mathbb F_2[[E]]^{\,n+1}.
 \tag{A.82.9}
\]
At every finite quotient \(E/U\), finite-length
Krull--Schmidt cancellation makes the reduction of
\(\overline R\) free of rank \((n+1)/2\).
For completeness, this cancellation uses the finite-length
Fitting decomposition: an endomorphism of an indecomposable
module is invertible or nilpotent, so its endomorphism ring
is local; the usual summand-exchange argument then gives
unique indecomposable multiplicities. Doubling those
multiplicities in (A.82.9) permits cancellation.

The finite sets of bases have the inverse-limit finite
intersection property: a basis at a common finer quotient
induces compatible bases at any finite collection of levels.
No arbitrary prechosen bases are being retained. Compactness
gives \(\overline R\simeq\mathbb F_2[[E]]^{(n+1)/2}\).
Lift such a basis to \(R_2\). The cokernel of the resulting
map from \(\Lambda_E^{(n+1)/2}\) is a compact pro-two
abelian group C with \(C=2C\), hence is zero by its finite
quotients. The surjection splits by projectivity; its kernel
has zero reduction modulo two and vanishes by the same
argument. This proves (A.82.3).
\end{proof}

\emph{The actual incidence map still has to be invertible.}\par
In the three-class case let \(e_i\) be the identity-coset
vector of \(\mathbb Z_2[[E/C_i]]\). The particular map
\[
 d:\mathbb Z_2[[E]]^2\longrightarrow R_2,
       \qquad d(a,b)=a(e_2-e_1)+b(e_3-e_1)
 \tag{A.82.10}
\]
has cokernel
\[
 \operatorname{coker}d\simeq
 \ker\bigl(\mathbb Z_2[[E/H]]\longrightarrow\mathbb Z_2\bigr),
       \qquad H=\overline{\langle C_1,C_2,C_3\rangle}.
 \tag{A.82.11}
\]
Indeed identifying the three identity vertices in the
permutation modules imposes exactly the relations of H;
this is also the component module of the two-edge coset
graph. The reduced component module vanishes exactly when
\(H=E\), since a proper closed H is contained in a proper
open subgroup and its nontrivial finite coset module detects
the difference. For the one-class case fix an actual \(x\in F\),
put \(H_x=\overline{\langle C_1,x\rangle}\), and
write the entire one-loop map as
\[
 \begin{gathered}
 d_x:\mathbb Z_2[[E]]\longrightarrow R_2,
       \qquad d_x(a)=a(xe_1-e_1),\\
 \operatorname{coker}d_x\simeq
 \ker\bigl(\mathbb Z_2[[E/H_x]]\longrightarrow\mathbb Z_2\bigr).
 \end{gathered}
 \tag{A.82.12}
\]
Indeed the relations \(ge_1=gxe_1\), together with the
existing \(gce_1=ge_1\) for \(c\in C_1\), identify
exactly the cosets \(gH_x\) of E. This proves the
cokernel description directly and makes surjectivity
equivalent to \(H_x=E\).

In either case a surjective displayed incidence map is
an isomorphism: after choosing the target basis from
(A.82.3), it is a surjective endomorphism of a finite-rank
free completed group-ring module. Its reductions over
\((\mathbb Z/2^a)[E/U]\) are surjections of finite
sets of equal order, hence injective; these quotients
separate the original module. Thus target freeness
leaves the selected incidence map's surjectivity---the
required generation assertion---unproved.
The full completed group algebra is not local; a unit test
after augmentation or a pro-two specialization does not
supply that assertion. No marked peripheral conjugator is
claimed.
